\documentclass[10pt,a4paper]{amsart}
\usepackage[margin=19mm]{geometry}
\usepackage{amsmath,amssymb,mathtools}
\usepackage[scr=rsfs,cal=euler]{mathalfa}
\usepackage{xfrac}
\usepackage[unicode,hidelinks,bookmarks=false]{hyperref}
\newcommand{\ie}{i.e.\ }
\newcommand{\cf}{cf.\ }
\newcommand{\resp}{respectively\ }
\newcommand{\Subsection}{\subsection}
\providecommand{\nobreakdash}{\nobreak\hbox{-}}
\setlength{\emergencystretch}{2em}
\multlinegap0pt
\usepackage{amssymb}
\newtheorem{theorem}{Theorem}
\title{Blowup analysis of a Camassa-Holm type equation with time-varying dissipation}
\author{Yonghui Zhou, Xiaowan Li, Shuguan Ji}
\date{Source: arXiv:2603.26534v1 (CC BY 4.0)}
\begin{document}
\maketitle

\noindent\emph{Source and licence.} Blowup analysis of a Camassa-Holm type equation with time-varying dissipation. Version arXiv:2603.26534v1 (CC BY 4.0), distributed by arXiv under the Creative Commons Attribution 4.0 International licence, http://creativecommons.org/licenses/by/4.0/, as recorded in the arXiv licence metadata for this exact version and shown on the abstract page https://arxiv.org/abs/2603.26534v1. Full licence notice: \texttt{licenses/arxiv-2603.26534-CC-BY-4.0.txt}.\par\smallskip

\noindent\textit{Editorial scope.} Counted unit: the article's theorem the206 (source0.tex lines 655--661). The article's complete original proof of it is delivered with it (source0.tex lines 662--772, one \texttt{proof} environment of 111 lines). No mathematics is written, repaired or re-derived: the statement and the proof are the article's own lines, in the article's own notation, and no retained line is substituted or re-flowed.  Retained uncounted as the source-local context the proof needs: lines 190--193; lines 553--562; lines 582--586.  Nothing was omitted from any retained span, so the package carries no omission record.  No retained line contains a citation, so the wrapper carries no bibliography and no bibliography entry.  No retained line contains a figure environment, an image-inclusion command or drawing code, so no image is delivered and the package declares no asset dependency.  Editorial formatting only, in addition: an AMS \texttt{amsart} preamble that restates the article's own declarations and macros used by the retained lines, namely the environments \texttt{theorem} on separate counters, together with the standard packages those lines need.  The retained environments print in this document as Theorem 1 in order of appearance, whereas the article prints the same items with the numbering of its own full document.  The complete native source of the article as distributed in the arXiv e-print for this exact version is bundled unchanged as \texttt{sources/197290/source0.tex}.
\begin{equation}
    u_t - u_{txx} + 3u^2 u_x + \lambda(t)(u - u_{xx}) = u u_{xxx} + 2 u_x u_{xx},
    \label{103}
\end{equation}

\begin{theorem}
\label{the205} Let $u_{0}\in H^{s} (s>\frac{3}{2})$ be given and $T$
be the maximal existence time of the corresponding solution $u(t,x)$ to problem \eqref{103}.
Assume that $\lambda(t)\leq \delta$ and some $x_{0}\in \mathbb{R}$ such that
\begin{equation}
u_{0}'(x_{0})<-\delta-\sqrt{\delta^{2}+2K},
\label{2015}
\end{equation}
where $K=\frac{\sqrt{2}}{2}\|u_{0}\|^{3}_{1}+\frac{5}{2}\|u_{0}\|^{2}_{1}.$ Then the corresponding solution of the Cauchy problem \eqref{103} blows up in finite time.
\end{theorem}

\begin{equation}
m'(t)=-\frac{1}{2}m^{2}-\lambda(t) m +u^{2}-P\ast\left(u^{2}+\frac{u_{x}^{2}}{2}\right)
-P\ast h(u)+h(u),\ \ a.e.\ on\ (0,T).
\label{2018}
\end{equation}

\begin{theorem}
\label{the206}
Under the conditions of Theorem \ref{the205}. If $T_{1}^{\ast}$ is finite, then
\begin{equation*}
\liminf_{t\rightarrow T_{1}^{\ast}}\inf_{x\in \mathbb{R}}u_{x}(t,x)(T_{1}^{\ast}-t)=-2.
\end{equation*}
\end{theorem}

\begin{proof}
The argument of proof is divided into six steps.

\noindent\textbf{Step 1. Key differential inequality.}
From \eqref{2018}, we obtain that $m(t)$ satisfies
\begin{align}
m'(t) +\frac{1}{2} \big( m(t) + \lambda(t) \big)^2 = H(t),
\label{2022}
\end{align}
where $$H(t)=u^{2}-P\ast\left(u^{2}+\frac{u_{x}^{2}}{2}\right)
-P\ast h(u)+h(u)+\frac{1}{2}\lambda^{2}(t).$$ It is easy check that $H(t)$ is uniformly bounded: $|H(t)| \leq K+\frac{\delta^{2}}{2}:=K_{1}$ for all $t \in [0, T_{1}^*)$.

\noindent\textbf{Step 2. Introduction of auxiliary variables.}
Let $$z(t) = m(t) + \lambda(t).$$ Then $z(t) \to -\infty$ as $t \to T_{1}^*$. By the continuity of $\lambda(t)$, the limit $\bar\lambda = \lim_{t \to T_{1}^*} \lambda(t)$ exists. For any $\varepsilon > 0$, there exists $\delta > 0$ such that $$|\lambda(t) - \bar\lambda| < \varepsilon,\text{ for }T_{1}^* - t < \delta.$$

Define $\theta(t) = m(t) + \bar\lambda$. Then $\theta(t) \to -\infty$, and $z(t) = \theta(t) + (\lambda(t) - \bar\lambda)$.

\noindent\textbf{Step 3. Estimating the difference $z^2(t) - \theta^2(t)$.}
We compute
\begin{equation}
|z^2(t) - \theta^2(t)| = |2\theta(t)(\lambda(t) - \bar\lambda) + (\lambda(t) - \bar{\lambda})^2|
\leq 2\varepsilon |\theta(t)| + \varepsilon^2.
\label{2023}
\end{equation}


\noindent\textbf{Step 4. Reformulating the equation for $\theta(t)$.}
From \eqref{2022} and $m'(t) = \theta'(t)$, we have
$\theta'(t) = -\frac{1}{2} z^2(t) + H(t).$
Using \eqref{2023}, we obtain
\begin{align}
\theta'(t) &\leq -\frac{1}{2}\theta^2(t) + \varepsilon |\theta(t)| + \frac{1}{2} \varepsilon^2 + K_{1},\nonumber\\
\theta'(t) &\geq -\frac{1}{2}\theta^2(t) - \varepsilon |\theta(t)| -\frac{1}{2} \varepsilon^2 - K_{1}.
\label{2024}
\end{align}
Since $\theta(t) < 0$, let $\phi(t) = -\theta(t) > 0$. Then $\phi(t) \to +\infty$ and \eqref{2024} becomes
\begin{align}
\phi'(t) &\geq \frac{1}{2} \phi^2(t) - \varepsilon \phi(t) - K_2,\nonumber\\
\phi'(t) &\leq \frac{1}{2}\phi^2(t) + \varepsilon \phi(t) + K_2,
\label{2025}
\end{align}
where $K_2 = K_1 + \frac{1}{2} \varepsilon^2$.

%\noindent\textbf{Step 5. Dominant balance and bounds for $\phi'(t)$.}
%Since $\phi(t) \to +\infty$, we can choose $t_1 < T^*$ such that for all $t \in [t_1, T^*)$, the following inequality holds:
%$$\frac{1}{4} \phi^2 (t)\geq \varepsilon \phi(t) + M_1.
%$$
%This is possible because the left-hand side grows quadratically while the right-hand side grows linearly in $\theta$. Under condition (5.1), we analyze the inequalities in (5).

%From the lower bound in (5):

%\begin{equation}
%\phi'(t) \geq \frac{1}{2} \phi^2(t) - \varepsilon \phi(t) - M_1
%=\frac{1}{4} \phi^2(t) + \left( \frac{1}{4} \phi^2(t) - \varepsilon \phi(t) - M_1 \right)
%\geq \frac{1}{4}\phi(t)^2 .
%\end{equation}


%From the upper bound in (5):

%\begin{equation}
%\phi'(t) \leq \frac{1}{2} \phi^2(t) + \varepsilon \phi(t) + M_1
%=\frac{3}{4} \phi^2(t) - \left(\frac{1}{4} \phi^2(t) - \varepsilon \phi(t) - M_1 \right)
%\leq \frac{3}{4} \phi(t)^2.
%\end{equation}


%Combining (5.2) and (5.3) yields:
%\[
%\frac{1}{4} \phi^2(t) \leq \phi'(t) \leq \frac{3}{4} \phi^2(t). \tag{6}
%\]

\noindent\textbf{Step 5. Precise blow-up rate for $\phi(t)$.}
Dividing \eqref{2025} by $\phi^2(t)$ gives
\begin{align}
\frac{1}{2} - \frac{\varepsilon}{\phi(t)} - \frac{K_2}{\phi^2(t)} \leq \frac{\phi'(t)}{\phi^2(t)} \leq \frac{1}{2} + \frac{\varepsilon}{\phi(t)} + \frac{K_2}{\phi^2(t)}.
\label{2026}
\end{align}
Since $\phi(t) \to +\infty$, for any $\varepsilon' > 0$, there exists $t_2 \in [t_1, T_{1}^*)$ such that for $t \in [t_2, T_{1}^*)$,
\begin{align}
\frac{1}{2} - \varepsilon' \leq \frac{\phi'(t)}{\phi^2(t)} \leq \frac{1}{2} + \varepsilon'.
\label{2027}
\end{align}
Note that $\frac{\phi'(t)}{\phi^2(t)} = -\frac{d}{dt} \left( \frac{1}{\phi(t)} \right)$. Integrating \eqref{2027} with respect to $t$ from $t$ to $T_{1}^*$ yields
$$ \left( \frac{1}{2} - \varepsilon' \right)(T_{1}^* - t) \leq \frac{1}{\phi(t)} \leq \left( \frac{1}{2} + \varepsilon' \right)(T_{1}^* - t).
$$
Hence,
\begin{align}
\lim_{t \to T_{1}^*} \phi(t)(T_{1}^* - t) = 2.
\label{2028}
\end{align}

\noindent\textbf{Step 6. Returning to $m(t)$.}
Since $\phi(t) = -m(t) - \bar\lambda$, we have
\begin{align*}
-m(t) - \bar\lambda \sim \frac{2}{T_{1}^* - t},
\end{align*}
which implies
\begin{align*}
m(t)(T_{1}^* - t) \sim -2 - \bar\lambda (T_{1}^* - t) \to -2.
\end{align*}
Therefore,
\begin{align*}
\liminf_{t \to T_{1}^*} m(t)(T_{1}^* - t) = -2,
\end{align*}
i.e.,
\begin{align*}
\liminf_{t \to T_{1}^*} \inf_{x\in \mathbb{R}}u_{x}(t) \cdot (T_{1}^* - t) = -2.
\end{align*}
This completes the proof.
\end{proof}

\end{document}
